% !TeX root = ../main.tex
% Risk, contracts and project delivery.
\chapter{Risk, Contracts and Project Delivery}
\label{ch:delivery}

Ground risk cannot be removed by investigation, and the commercial question is
therefore how it is allocated. This chapter closes the book by returning to
the argument with which it opened. Most disputes in underground construction
turn on two things, what was known about the ground and which party bears the
consequences of the remaining uncertainty. The contractual instruments examined
here make that allocation explicit; they do not eliminate the
uncertainty.

\section{Ground risk and baselines}
\label{sec:ground-risk-baselines}

\subsection{The nature of ground risk}
\label{sec:nature-ground-risk}

Ground risk has three properties that distinguish it from most construction
risk. It cannot be eliminated by investigation, only reduced and
characterised. It is unevenly distributed along an alignment, concentrated at
features that may not have been found. Finally, the party best placed to manage it
during construction (the contractor) is not always the party best placed
to have investigated it.

The second of those properties has a consequence for how exposure is
presented. Figure~\ref{fig:exposure-by-zone} aggregates predicted
third-party exposure along an illustrative alignment by the zone each building
falls in,
which is the form in which it can be owned: a register entry needs a location,
a responsible party and a mitigation, and a corridor-wide average provides none
of them.

\begin{figure}[htbp]
\centering
\datafigure[0.78\textwidth]{ch28/receptor-risk}
\caption[Third-party exposure aggregated by zone]{%
Third-party exposure aggregated by zone: the share of buildings
within 800\,m of each zone centre predicted to settle by more than 10\,mm along
an illustrative twin-tunnel alignment. The figures beside each bar are the
buildings affected and the buildings assessed. Exposure is normally
concentrated at particular zones rather than spread evenly along the alignment.
Synthetic demonstrator data (Appendix~\ref{app:software}).}
\label{fig:exposure-by-zone}
\end{figure}

Risk management runs on a register: hazards identified systematically,
assessed for likelihood and consequence, assigned an owner and a mitigation,
and reviewed as the project progresses. Its value depends on being kept
current. A register compiled to satisfy a requirement at design stage and never
revisited records compliance only; one updated as the ground is revealed
allows the project to act on what it finds.

Quantitative assessment is justified where the exposure warrants it, with
schedule and cost risk analysis producing a distribution of outcomes rather
than a point estimate. The insurance-driven code of practice for the risk
management of tunnel works~\citep{itig2012code}, developed after a sequence of
major losses, made structured risk management a contractual requirement
rather than a matter of good practice, and it places obligations on client,
designer and contractor alike.

\subsection{The geotechnical baseline report}
\label{sec:baseline-report}

A geotechnical baseline report~\citep{essex2007gbr} states the ground
conditions the contract assumes. It is neither a factual report, which presents data, nor an
interpretative report, which presents an opinion: it is a contractual
statement that conditions no worse than those described will be treated as
having been allowed for, and conditions worse than described will be treated
as a changed condition.

The requirement in writing one is testability. Every
baseline statement must be capable of being compared against what is
encountered, by an agreed method, during construction. ``Occasional fault
zones'' is not a baseline; ``not more than 200\,m of ground classified as Class
V by the stated method'' is. An ambiguous statement leads to a dispute
rather than avoiding one.

\subsection{Differing conditions and claims}
\label{sec:differing-conditions}

Where the baseline is testable, the mechanism works as intended. The contractor prices the baselined
conditions without contingency for the unknown, the client carries the risk of
conditions being worse (Figure~\ref{fig:baseline-mechanism}), and the comparison is made against a stated method
rather than argued from first principles. Where the baseline is silent or
untestable, the parties fall back on argument about what a reasonable
contractor should have anticipated, and such arguments are how underground
disputes become protracted.

\begin{figure}[htbp]
\centering
\datafigure[0.62\textwidth]{ch28/baseline-mechanism}
\caption{A testable baseline applied, with illustrative values: the cumulative length of Class~V ground encountered along a drive against a baseline of 200~m. Ground up to the baseline is priced by the contractor; the excess is a changed condition.}
\label{fig:baseline-mechanism}
\end{figure}

\section{Approvals and performance requirements}
\label{sec:approvals-performance}

A major tunnel is approved before it is designed in detail, and the approval
fixes the outcomes the design must deliver. The criteria a designer works to
on settlement, vibration, ground-borne noise and groundwater are therefore
set, in large part, by an environmental assessment conducted years earlier and
written into the instruments that authorise the works. Reading those
instruments correctly, and translating each requirement into a design
quantity, is a design task rather than an administrative one.

\subsection{How the requirements arise}
\label{sec:requirements-arise}

In Victoria a project likely to have significant effects is assessed through
an Environment Effects Statement under the Environment Effects Act
1978~\citep{vic1978eeact}. The Minister for Planning sets scoping
requirements, the proponent prepares the statement with its technical studies,
it is exhibited and usually examined by an inquiry, and the Minister then issues
an assessment that the decision-makers must consider but that does not itself
bind them~\citep{dtp2024ees}. The binding element is the environmental
management framework that accompanies the approval. Its environmental
performance requirements define the outcomes that must be achieved during the
further design, construction and operation; they are given force through the
planning scheme provisions that authorise the project and through the delivery
contracts, and compliance is checked by independent
audit~\citep{wda2017ees}. In New South Wales the equivalent instrument is the
set of conditions attached to an approval of critical State significant
infrastructure under the Environmental Planning and Assessment Act
1979~\citep{nsw1979epa,nsw2020m4m5coa}.

The two systems write their requirements differently, and the difference
matters to the designer. A Victorian requirement for ground movement typically
prescribes a process rather than a number: prepare a geotechnical model, use
it to predict settlement and strain, survey the condition of buildings and
utilities, agree acceptable criteria with each owner, monitor, and repair or
compensate any damage~\citep{wda2017ees,dtp2022srlassessment}. The numerical
criterion is then set during design, by the designer, and is part of what is
reviewed. A New South Wales condition is more often numerical and is fixed in
the approval itself. Vibration and ground-borne noise are numerical in both,
expressed as targets referenced to the standards of
Chapter~\ref{ch:construction}.

\subsection{From requirement to design quantity}
\label{sec:requirement-to-design}

Table~\ref{tab:requirements-design} traces published requirements to the
quantity each one constrains and to the method of this book that evaluates it.
Three translations show the pattern.

\emph{Settlement.} One published condition limits a sensitive building to a
settlement of 20~mm, an angular distortion of 1 in 500 and a limiting tensile
strain of 0.1~per~cent~\citep{nsw2020m4m5coa}. The Gaussian trough of
Section~\ref{sec:gaussian-trough} converts the settlement limit into a limit on
volume loss. For the example of that section (a 6~m tunnel at 20~m depth
with $K = 0.5$, so that $i = 10$~m), a volume loss of 1~per~cent gives
$S_{\max} = 11.3$~mm, and the limit of 20~mm therefore admits
$V_L \leq 20/11.3 = 1.8$~per~cent (Figure~\ref{fig:settlement-volume-loss}). The steepest slope of that trough,
$S_{\max} e^{-1/2}/i = 20 \times 0.607/10\,000 = 1.2 \times 10^{-3}$, or
1 in 820, is inside the angular distortion limit, so settlement governs here,
and the strain criterion is checked by the damage assessment of
Section~\ref{sec:damage-assessment}. The volume loss in turn sets the face
pressure and annulus grouting of Chapter~\ref{ch:tbm}, and the trigger levels of
Section~\ref{sec:triggers-interpretation} are placed below the limit by the
margin the prediction uncertainty requires.

\begin{figure}[htbp]
\centering
\datafigure[\textwidth]{ch28/settlement-volume-loss}
\caption{The settlement condition as a limit on volume loss, for a 6~m tunnel at 20~m depth with $K = 0.5$. Left: $S_{\max}$ reaches 20~mm at $V_L = 1.8$~per~cent. Right: the steepest slope there is $1.2 \times 10^{-3}$, inside 1 in 500, so settlement governs.}
\label{fig:settlement-volume-loss}
\end{figure}

\emph{Groundwater.} The same approval limits the inflow to an operating tunnel
to 1~L/s in any kilometre, not averaged over the drive~\citep{nsw2020m4m5coa},
that is $10^{-6}$~m$^3$/s per metre. The steady inflow to a drained tunnel of
Section~\ref{sec:inflow-prediction}, $q = 2\pi k_w H/\ln(2H/r_0)$, gives for a
head of $H = 20$~m above a tunnel of radius $r_0 = 3$~m
$q/k_w = 2\pi \times 20/\ln(13.3) = 48.5$~m, so that the condition is met
without a membrane only where the rock mass permeability is below
$10^{-6}/48.5 \approx 2 \times 10^{-8}$~m/s (Figure~\ref{fig:inflow-permeability}). Above that value the requirement
decides the choice of Section~\ref{sec:drained-vs-tanked} in favour of a tanked
lining or of systematic grouting, whatever the structural or carbon arguments.
A Victorian requirement for the same quantity is an outcome: minimise the
change in groundwater level and the inflow~\citep{wda2017ees}. It leaves
the threshold to be argued in design.

\begin{figure}[htbp]
\centering
\datafigure[0.62\textwidth]{ch28/inflow-permeability}
\caption{Steady inflow per kilometre to a drained tunnel of radius 3~m against rock mass permeability, for heads $H$ above the axis. The condition of 1~L/s in any kilometre is met without a membrane below $k_w = 2 \times 10^{-8}$~m/s at $H = 20$~m.}
\label{fig:inflow-permeability}
\end{figure}

\emph{Vibration and ground-borne noise.} The Victorian targets of one published
framework take human comfort from the vibration dose values of BS~6472-1,
0.2 to 0.4~m/s$^{1.75}$ by day and 0.1 to 0.2~m/s$^{1.75}$ at night for a
residence; structural damage from DIN~4150-3; blasting from AS~2187.2, at
5~mm/s for 95~per~cent of blasts and 10~mm/s at most at a sensitive site; and
ground-borne noise from an internal level of 40~dB(A) in the evening and
35~dB(A) at night as $L_{\mathrm{Aeq},15\,\mathrm{min}}$~\citep{wda2017ees}.
A New South Wales approval sets the same ground-borne noise
levels~\citep{nsw2020m4m5coa}. Each is converted into a stand-off distance, a
charge per delay or a restriction on night working by the relations of
Section~\ref{sec:vibration-environment}.

\begin{table}[htbp]
\caption{Published performance requirements and the design quantities they
constrain. The values are those of the instruments cited and apply to those
projects; a new project takes its own.}
\label{tab:requirements-design}
\centering
\small
\begin{tabularx}{\textwidth}{>{\raggedright\arraybackslash}p{0.30\textwidth} >{\raggedright\arraybackslash}X >{\raggedright\arraybackslash}p{0.22\textwidth}}
\toprule
\textbf{Requirement} & \textbf{Design quantity it constrains} & \textbf{Method} \\
\midrule
Settlement of a sensitive building $\leq 20$~mm, distortion $\leq$ 1 in 500, tensile strain $\leq 0.1$\,\% \citep{nsw2020m4m5coa} & Volume loss, hence face pressure and annulus grouting & Sections~\ref{sec:gaussian-trough} and~\ref{sec:damage-assessment} \\
Settlement criteria agreed with each owner after condition survey \citep{wda2017ees} & The criterion itself, set and justified in design & Section~\ref{sec:damage-assessment} \\
Operational inflow $\leq 1$~L/s in any kilometre \citep{nsw2020m4m5coa} & Drained or tanked lining; extent of grouting & Sections~\ref{sec:inflow-prediction} and~\ref{sec:drained-vs-tanked} \\
Vibration dose value, residence, 0.2--0.4~m/s$^{1.75}$ by day \citep{wda2017ees} & Stand-off distance and machine operation & Section~\ref{sec:vibration-environment} \\
Ground-borne noise $L_{\mathrm{Aeq},15\,\mathrm{min}}$ 40~dB(A) evening, 35~dB(A) night \citep{wda2017ees,nsw2020m4m5coa} & Night working; respite and mitigation & Section~\ref{sec:vibration-environment} \\
Blast vibration 5~mm/s for 95\,\% of blasts, 10~mm/s maximum, sensitive site \citep{wda2017ees} & Charge per delay at each distance & Section~\ref{sec:vibration-environment} \\
\bottomrule
\end{tabularx}
\end{table}

Two cautions apply to every row. First, a requirement is a limit on the outcome, not
a design value: the design must reach it with the prediction uncertainty of
Chapter~\ref{ch:integrated-design} accounted for, which places the design
target and the trigger levels below the requirement rather than at it. Second, a
requirement belongs to the approval in which it was written; values from
another project are evidence of what regulators have accepted, not criteria
that transfer.

\section{Contracts, cost and delivery}
\label{sec:contracts-cost-delivery}

Contract models differ mainly in where they place ground risk. Design-bid-build
leaves the client with design risk and the contractor with construction risk,
and the interface between them is where ground risk lands. Design and construct
transfers more to the contractor, who prices the uncertainty and, where it
cannot be quantified, prices it high or declines to bid. Early contractor
involvement brings construction expertise into design, where it can influence
the decisions that determine constructability. Alliance and other collaborative
models share risk explicitly and align commercial interests, at the cost of
requiring a client organisation able to participate rather than only
administer.

No model eliminates ground risk; each decides who is paid to carry it. The
failure mode common to all is transferring unquantifiable risk to a party
without the information to price it, which produces either an inflated
contingency the client pays regardless, or a contractor who prices optimistically
and litigates later.

Procurement determines how much of this is settled before the bid. Tender
period investigation, reference designs, and a properly written baseline all
reduce the uncertainty a bidder must price. Dispute avoidance mechanisms
(standing boards that follow the project and address issues as they arise)
resolve matters while the evidence is fresh and the relationship is intact, and
their record is better than that of adjudication after the event.

Cost and programme estimation for underground works must be probabilistic,
because the distribution of outcomes is skewed: the best case is bounded by
the theoretical cycle, and the worst case is not. Reference class forecasting
(comparing against the outturn distribution of similar completed projects
rather than building up from first principles) outperforms bottom-up
estimation where a reference class exists, because it captures the
difficulties that bottom-up estimates omit by construction. Advance rates should come from
comparable projects by ground class, with realistic utilisation
(Chapter~\ref{ch:tbm}), and contingency should be set from the risk analysis
rather than as a percentage. The risk analysis yields a distribution of cost
built from the register; its central value is the estimate, an upper
percentile is the figure funded, and the difference between them is the risk
allowance (Figure~\ref{fig:cost-distribution}). A percentage uplift cannot produce this quantity, because it is not
derived from the risks it is meant to cover.

\begin{figure}[htbp]
\centering
\datafigure[\textwidth]{ch28/cost-distribution}
\caption{An illustrative outturn cost distribution, bounded below and skewed to the right. Left: the density, with the tail above P90 shaded. Right: the cumulative curve, on which the estimate is the P50, the funded figure the P90, and the risk allowance the difference.}
\label{fig:cost-distribution}
\end{figure}

Safety in design
obligations require that risk be designed out where practicable, which in
underground work means alignment and method decisions taken with construction and
maintenance safety as explicit criteria. The most effective safety measures
are the earliest ones: an alignment that avoids the fault zone, a method that
keeps people away from an unsupported face, and a design that can be inspected
and maintained without exposing anyone to the conditions that made it
difficult to build.

\section{Digital delivery and model liability}
\label{sec:digital-delivery}

The baseline report of Section~\ref{sec:baseline-report} is a document, and the contractual
machinery around ground risk was built for documents. Projects are now
delivered with models: a ground model, an asset model, and increasingly an
operating twin. The commercial framework has not caught up, and the gap shows
in three questions that arise on every digitally delivered project and are
answered on very few: what was promised, who owns the result, and who
carries the consequence when a decision taken from a model proves to have
been taken from a wrong one.

\subsection{The model as a contract deliverable}
\label{sec:model-deliverable}

A model becomes a deliverable only if the contract says what it must contain,
and it is more useful to specify the information required than
the software used. The ISO~19650 series~\citep{iso19650} provides the
structure: the client states its information requirements, the delivery team
responds with a plan for meeting them, and information is produced,
checked and issued through a common data environment with defined states of
approval.

Three specifications matter most. The first is purpose: what decisions the
model must support, since a model adequate for clash detection is not adequate
for settlement assessment and neither is adequate for asset management.
The second is format: an open, documented exchange format, because a deliverable readable
only under the licence that produced it leaves the client dependent on that licence.
The third is fidelity: the level of detail and, more importantly, the
reliability of the information, that is, which parts are measured, which
interpreted and which assumed.

The third specification connects the model to the baseline. A ground model
carrying its uncertainty explicitly, as
Chapter~\ref{ch:ground-modelling} describes, can support a baseline that
states expected conditions and their variation rather than a single line on a
long-section. This describes the ground more faithfully than the conventional
document and is also more testable, which works in both directions
commercially: it is harder to argue
that conditions were unforeseen when the model published the range in advance,
and harder to argue that they were foreseen when it did not.

\subsection{Ownership, warranty and handover}
\label{sec:model-ownership}

Ownership of a model is rarely addressed and rarely simple. The ground model
is typically built by the designer from data the client procured, updated by
the contractor during construction, and needed by the operator for sixty years
after both have left. Intellectual property in the model, in the software
configuration and in the underlying data may sit with different parties, and a
handover that transfers the file without transferring the right to use, modify
and sub-licence it delivers an object the operator cannot lawfully maintain.

Warranty requires equal care. A designer will warrant reasonable skill and
care in preparing a model; no competent party will warrant that a ground model
is correct, because it is an interpretation of sparse data. What can sensibly
be warranted is that the model represents the data
available, that its assumptions are stated, and that it has been prepared to
the agreed standard. Contracts that attempt more than this are unlikely to be
accepted by a party that has understood them.

Handover should transfer four things, and the omission of the last is the
usual failure: the model itself, in the agreed open format; the data behind
it, so that it can be rebuilt rather than only opened; the record of what
was assumed and by whom; and a named party responsible for keeping it current.
Without that party the model is accurate on the day of handover and steadily less
so thereafter
(Chapter~\ref{ch:assets}).

\subsection{Liability for model-informed decisions}
\label{sec:model-liability}

The exposure that concerns practitioners most is the one least settled: what
happens when a decision is taken with the assistance of a model, or of a
statistical prediction, and the outcome is bad.

The principles that apply are not new. A professional who relies on a tool
remains responsible for the reliance being reasonable, which requires knowing
the tool's limits. A supplier who provides a model for a stated purpose may be
liable if it is unfit for that purpose, which is why the stated purpose in
Section~\ref{sec:model-deliverable} matters commercially as well as
technically. A party who uses output beyond the purpose for which it was
supplied has generally assumed the risk of doing so.

The new difficulty is evidential. When a decision is informed by a
statistical model that has since been retrained, the model that produced the
advice no longer exists, and reconstructing the basis of the decision may be
impossible. This is the practical reason for the record-keeping required in
Chapter~\ref{ch:digital-twin}: versioned models, logged predictions, and a
retraining trail. The record protects the engineer as much as it assures
the client. It is cheap to establish at the outset and cannot be
created retrospectively.

The prudent position for a designer is therefore to state the purpose of
every model issued, to state what it does
not cover, to keep the version that supported each decision, and to avoid contractual
language implying that a model predicts conditions rather than represents an
interpretation of them. These measures do not obstruct digital delivery, and
they make it defensible.

\chapterreferences
